\chapter{Introdução}


\section{Contextualização}

A plataforma tecnológica BEM-TE-VI da Universidade de Brasília abriga o projeto v-FES cycling, que combina o ciclismo
com o apoio da estimulação elétrica funcional (FES) e a realidade virtual na reabilitação neurológica.
O acompanhamento dos participantes do projeto depende de avaliações regulares realizadas pela equipe de fisioterapia, e
cada avaliação utiliza uma ferramenta específica: uma versão adaptada da Classificação Internacional de Funcionalidade, Incapacidade e Saúde (CIF),
a escala ASIA/ISNCSCI, a Escala de Alcance de Metas (GAS), o Miniexame de Estado Mental (MEEM), o eletrodiagnóstico, a densitometria pelo método DEXA
e a anamnese.

Os dados obtidos por meio desses instrumentos têm natureza distinta. A CIF, administrada pela Organização Mundial da Saúde, organiza
as funções e estruturas do corpo, as atividades e a participação e os fatores ambientais em códigos hierárquicos, aos quais são atribuídos qualificadores
\cite{who2001icf}. O MEEM gera pontuações, o eletrodiagnóstico registra medições por grupos musculares,
a densitometria registra medições regionais e de composição corporal, e a anamnese armazena texto livre.
As ferramentas também diferem na estrutura dos campos e nas regras de preenchimento.

O plano de trabalho da iniciação científica associado ao projeto descreve os registros dessas avaliações
como distribuídos entre arquivos de texto, tabelas, arquivos PDF, imagens e formulários, sem um modelo padronizado. Esse
mesmo plano prevê a criação de um banco de dados clínico estruturado e validado, que deveria servir de base para um futuro
sistema de inteligência artificial generativa voltado para a personalização do plano de treinamento \cite{ochoadiaz2026plano}. 
desse sistema são objeto de um estudo de doutorado associado a este projeto.

Com o objetivo de reunir os registros em um único local, foi desenvolvido o aplicativo Bem-Te-Vi, que consiste
em uma aplicação Flutter, uma API FastAPI e um banco de dados PostgreSQL. A primeira versão incluía login, cadastro
de pacientes e formulários de avaliação. Nesta etapa, foi adicionado um módulo de avaliação CIF com catálogos
versionados obtidos a partir das tabelas da equipe, cálculo hierárquico de pontos, esboço, tabela resumida, histórico e
comparação de avaliações do mesmo paciente. Também foi revisada a implementação dos demais módulos, com a correção de erros pontuais,
entre eles a validação do cadastro, que agora exige nome, data de nascimento, número de identificação (CPF) e sexo, sendo que este último dado é necessário, pois
determina a variante da escala CIF utilizada \cite{higa2026bemtevi}.

Na versão atual, o aplicativo reúne todos os sete formulários, mas trata cada um deles separadamente, com estrutura própria e
regras de armazenamento independentes. Apenas o fluxo CIF oferece funções de acompanhamento, como histórico
e comparação. Os demais formulários registram cada avaliação separadamente.



\section{Problema e pergunta de pesquisa}
\label{sec:problema}

Os registros das avaliações do projeto v-FES cycling estão distribuídos entre arquivos de texto, planilhas, PDFs,
imagens e formulários sem um modelo padronizado \cite{ochoadiaz2026plano}. O aplicativo Bem-Te-Vi unificou parte desses registros em sete formulários clínicos,
cada um com sua própria estrutura e regras. Seu tamanho varia de
6 a 1.003 campos, conforme indicado na Tabela~\ref{tab:formularios}.

\begin{table}[htb]
	\centering
	\caption{Formulários clínicos do Bem-Te-Vi e quantidade de campos}
	\label{tab:formularios}
	\small
	\begin{tabular}{lrp{0.48\textwidth}}
		\toprule
		Formulário & Campos & Composição \\
		\midrule
		CIF & 985 / 1.003 & 984 respostas (sexo masculino) ou 1.002 (feminino) e a data da avaliação \\
		ASIA/ISNCSCI & 140 & 112 sensoriais, 20 motores e 8 complementares \\
		GAS & 134 & 13 metas com 10 campos cada e 4 de cabeçalho \\
		Eletrodiagnóstico & 88 & 12 grupos musculares com 7 campos cada e 4 de cabeçalho \\
		Densitometria (DEXA) & 42 & medidas regionais, tendências, composição corporal, peso e altura \\
		MEEM & 34 & 30 itens de pontuação e 4 de identificação e contexto \\
		Anamnese & 6 & 5 de identificação, contato e data e 1 de anotações \\
		\bottomrule
	\end{tabular}
	\fonte{Elaborada pelo autor a partir do código do aplicativo \cite{higa2026bemtevi}.}
\end{table}
  
No total, os formulários contêm 1.429 campos para a avaliação completa
de um paciente do sexo masculino e 1.447 para o sexo feminino, sendo que aproximadamente 69\% desses campos
pertencem à CIF, além disso, os instrumentos também diferem em sua organização interna. A GAS repete
o mesmo bloco de campos para cada objetivo, e o eletrodiagnóstico repete um bloco para cada
grupo muscular. A ASIA separa as respostas sensoriais das motoras, a CIF varia de acordo com o sexo
do paciente e a anamnese contém texto livre.
  
Concentrar esses dados em um só aplicativo não basta quando o significado e os vínculos de cada campo ficam implícitos
no código. Na versão atual, os formulários GAS, ASIA, anamnese, eletrodiagnóstico, MEEM e densitometria localizam o
paciente pelo nome completo e, se não o encontram, criam um novo cadastro apenas com esse nome. A tela da ASIA exibe um
campo de comentários que não é enviado ao servidor. O MEEM registra sempre o mesmo avaliador, "Dr. Examinador", e a
densitometria grava como data do exame a data do envio. Esses
comportamentos são esperados em formulários criados antes da definição de um modelo de dados comum e de regras acordadas
com a equipe, e por isso servem aqui como ponto de partida para os requisitos.
  
Esses casos correspondem a dimensões de qualidade descritas na literatura sobre registros eletrônicos de saúde. \citeonline{weiskopf2013dataquality} organizam a avaliação dessa qualidade em
completude, correção, concordância, plausibilidade e atualidade.
O comentário descartado compromete a completude, enquanto o avaliador fixo e a data de envio comprometem a correção. A
identificação pelo nome tem consequência mais grave. \citeonline{just2016patientmatching} mostram que divergências nos
campos de identificação estão entre as principais causas de cadastros duplicados de pacientes. No Bem-Te-Vi, um erro de digitação
no nome basta para dividir o histórico de um paciente entre dois cadastros.
  
O acompanhamento das avaliações depende dessa mesma base. Histórico, estado de cada avaliação, resumos e comparações só
são confiáveis quando cada resposta está ligada ao paciente, ao instrumento e à data corretos. Nenhum dos problemas
descritos se resolve apenas com código. É preciso definir como identificar o paciente, quem responde pela avaliação,
qual data vale como data do exame e quais campos devem ser acompanhados ao longo do tempo, e essas regras precisam
estar especificadas antes de qualquer reestruturação do banco ou do aplicativo. Há ainda a questão da troca de dados.
Uma base com estrutura própria resolve a organização interna, mas não conversa com outros sistemas de saúde. No
Brasil, a Rede Nacional de Dados em Saúde (RNDS) adota o padrão HL7 FHIR para o intercâmbio de informações clínicas,
e modelar a base segundo esse padrão desde o início evita uma conversão posterior. Por esses motivos, este trabalho
trata o problema como uma questão de engenharia de requisitos e parte da seguinte pergunta: quais requisitos e qual
modelo de dados devem orientar a evolução do aplicativo Bem-Te-Vi para estruturar os dados clínicos de forma
interoperável e apoiar o acompanhamento das avaliações?

A pergunta central se desdobra nas seguintes questões:

\begin{description}
	\item [P1] Os recursos do HL7 FHIR conseguem representar os sete formulários clínicos, de 6 a 1.003 campos,
	sem perda de informação? Quais estruturas, como blocos repetidos, variação por sexo, pontuações calculadas e
	texto livre, exigem tratamento específico?
	\item [P2] Quais requisitos funcionais, requisitos não funcionais e regras de negócio o aplicativo precisa
	atender para acompanhar as avaliações, tomando como base o que o módulo CIF já implementa?
	\item [P3] Qual é o nível de automação atual de cada escala, e quais etapas podem ser automatizadas por regra do
	próprio instrumento, sem depender de julgamento clínico?
\end{description}

\section{Justificativa}

Na reabilitação, a avaliação não é um registro isolado. \citeonline{rauch2008icf} descrevem o manejo clínico como um
ciclo de avaliar, definir metas, intervir e reavaliar, e defendem sistemas eletrônicos de documentação para sustentar
esse ciclo. No Brasil, a CIF foi adotada no Sistema Único de Saúde \cite{cns452:2012}, o que reforça a importância de
registrá-la de forma estruturada. No Bem-Te-Vi, apenas o fluxo CIF permite consultar o histórico e comparar avaliações.
Nos demais instrumentos, acompanhar um paciente ainda exige reunir registros manualmente. Definir o que centralizar e como relacionar os dados, antes da implementação, reduz o risco de retrabalho na base e no aplicativo.
A economia de tempo e de etapas manuais, porém, é uma expectativa, que só poderá ser medida depois da implementação.

Para a Engenharia de Software, o problema reúne atividades centrais da engenharia de requisitos: elicitar requisitos com
especialistas de outra área, modelar dados heterogêneos, especificar regras de negócio e manter a rastreabilidade entre
a fonte clínica, o requisito e o critério de aceitação. São atividades em que falhas são comuns. No levantamento
NaPiRE, requisitos incompletos e requisitos ambíguos aparecem entre os problemas mais frequentes relatados pela
indústria \cite{mendezfernandez2017napire}. Na saúde, \citeonline{cysneiros2002health} acrescenta exigências próprias
do domínio, como privacidade, confiabilidade e o uso de um vocabulário que a equipe de desenvolvimento não domina.
Para organizar o processo, este trabalho adota como referência a ISO/IEC/IEEE 29148 \cite{iso29148:2018}.

O trabalho também é viável no prazo da disciplina. O aplicativo, o módulo CIF, o protótipo no Figma e os formulários
já existem e podem ser analisados como artefatos. As regras de cálculo da CIF estão nas planilhas da equipe e no
catálogo do aplicativo, e as das demais escalas estão descritas na literatura de cada instrumento, o que permite
especificar e verificar o modelo sem depender de novas sessões com a equipe.



\section{Objetivos}

\subsection{Objetivo geral}

Especificar os requisitos e o modelo de dados interoperável, baseado no padrão HL7 FHIR, para estruturar os dados
clínicos do aplicativo Bem-Te-Vi e apoiar o acompanhamento das avaliações do projeto v-FES cycling, tomando como ponto
de partida o módulo de avaliação CIF desenvolvido pelo autor.

\subsection{Objetivos específicos}

\begin{enumerate}
	\item Inventariar os sete formulários clínicos do aplicativo, descrevendo a estrutura de campos de cada um, as funções
	já disponíveis e os problemas de registro encontrados no código, como a identificação do paciente pelo nome;
	\item Descrever o módulo de avaliação CIF implementado, com catálogo versionado, cálculo hierárquico das pontuações,
	rascunho, quadro resumo, histórico e comparação, e extrair dele os requisitos e as regras de negócio que já
	estão em funcionamento;
	\item Avaliar o nível de automação de cada escala, distinguindo os dados digitados, os valores que o aplicativo já
	calcula, os que podem ser calculados por regra definida no próprio instrumento e os que dependem de julgamento clínico;
	\item Especificar e priorizar os requisitos funcionais, os requisitos não funcionais e as regras de negócio, com fonte,
	justificativa e critério de aceitação para cada item;
	\item Modelar a base de dados clínicos com os recursos do HL7 FHIR, considerando os perfis adotados pela Rede
	Nacional de Dados em Saúde, e mapear os formulários para esse modelo, com o CIF em maior detalhe;
	\item Verificar o modelo quanto à cobertura dos campos dos sete formulários e quanto à resposta aos problemas de
	registro inventariados, usando instâncias de exemplo com dados fictícios;
	\item Consolidar a matriz de rastreabilidade e as pendências que orientarão a implementação da base no TCC 2.
\end{enumerate}

\section{Escopo e delimitação}

O trabalho se divide em duas etapas. O TCC 1 especifica os requisitos e o modelo de dados em HL7 FHIR e apresenta,
como resultado parcial, o módulo CIF já implementado. O TCC 2 implementa a base de dados e avalia os requisitos priorizados, em continuidade ao desenvolvimento previsto no plano de iniciação
científica \cite{ochoadiaz2026plano}. O escopo do TCC 2 será definido a partir dos resultados desta etapa. Além do
módulo CIF, já desenvolvido, nenhuma alteração no banco de dados ou no código do aplicativo faz parte do TCC 1. O
aplicativo e o protótipo entram como artefatos de análise.

Não há sessões de elicitação ou validação com a equipe de fisioterapia nesta etapa. As fontes dos requisitos são o
código do aplicativo, as planilhas da CIF que deram origem ao catálogo, o plano de trabalho e a literatura dos
instrumentos. A verificação do modelo usa instâncias de exemplo com dados fictícios, e as decisões que exigem
julgamento clínico ficam registradas como pendências para validação no TCC 2.

Os sete formulários clínicos são analisados pela estrutura, e não campo a campo. O modelo em HL7 FHIR deve cobrir
todos eles, mas o mapeamento detalhado se concentra na CIF, que já conta com histórico, quadro resumo e comparação
implementados e responde por cerca de 69\% dos campos. Os campos e funções que não forem tratados ficam registrados
como pendências para o TCC 2.

O plano de trabalho prevê que a base de dados alimente um sistema de inteligência artificial generativa. Os algoritmos
desse sistema são objeto do doutorado vinculado ao projeto. No TCC 1, cabe apenas descrever os dados que esse
módulo deverá consumir. Também ficam fora do escopo a interpretação clínica dos resultados, que continua com a equipe, a
integração efetiva com a RNDS, que exige credenciamento institucional, e outros padrões, como o openEHR, citados apenas
como referência.

Todas as atividades usam somente dados fictícios. O uso de dados reais de pacientes depende de autorização
institucional e não está previsto nesta etapa.

\section{Organização do trabalho}

%
